% s-248-p248-m003; original PDF 248, 249, 250, 251
\Aphorism{s-248-p248-m003}{347}
Gegen die modern seelenhafte Dichtung, als deren Grossmeister immer noch der gewaltige Ibsen dasteht, habe ich einen Einwand zu erheben, der bei Freund und Feind sonst eher als Ruhmestitel gilt: diese Dichtung ist von einer unerlaubten Zartheit und Zärtlichkeit in der Darstellung des modernen Culturmenschen. Es werden uns Individuen vorgeführt von einer Hautempfindlichkeit und inneren Verletzlichkeit, einer Verwöhntheit seelischer Ansprüche, einer messerscharf zugespitzten Singularität der Lebensbedingungen, dass man sich selbst dagegen wie ein vierschrötiger Allesverdauer und Glasscherben-Fresser vorkommt. Ich zweifle nicht, dass es da und dort solche mimosenhaften Seelchen gebe, und wenn das Leben in einer spöttisch-grossmüthigen Laune sie mit Glacéhandschuhen anfasst, mögen sie sich sogar eine Weile erhalten und fortpflanzen können: aber ausschliesslich mit derlei hyperästhetischen Nervenmenschlein die moderne Bühne bevölkern heisst uns doch der Nachwelt gegenüber blossstellen. Man lese einen Dialog solcher Sensitiven beiderlei Geschlechts, etwa zwischen Hilde und dem Baumeister, oder zwischen den Eltern des kleinen Eyolf, und alles Soldatische (im feinsten und gröbsten Sinne) wird sich Einem im Leibe empören bei dieser grenzenlosen Schonung und Verweichlichung, in die sich hier die Seelen einbalsamiren. Ob Zwei einander „verstehen“, ob sie halbe oder ganze, vorläufige oder entscheidende Proben dieses Verständnisses gegeben haben oder geben sollen, ob die Wage gegenseitiger Verbindlichkeit auf dieser oder jener Seite um einen Scrupel zuviel belastet sei, ob ihr Zusammenleben eine „wahre“ Ehe gewesen sei oder noch werden könne — das sind die Cardinalfragen dieser Dramatik, zu deren Lösung ein Raffinement von Spürsinn, Dialectik und Eiertanz aufgeboten wird, wie es bei wissenschaftlichen oder diplomatischen Fragen delicatester Art nicht gelegener käme: eine unanständige Mikroskopie und Beschnüffelung der Seele mit allen quarante manières geschulter Lüsternheit. Nun frage ich: welcher Mensch hat das Recht, für seinen Gallertquark von Seele einen derartigen Aufwand von Rücksichtnahme und behutsamster Scheu zu verlangen? Wer hat das Recht, seine hysterischen Eigenheiten als unübertretbare Existenzbedingungen aufstellen und sie dem einfachen Gange des Lebens gegenüber à outrance durchsetzen zu wollen? Wer darf in Ehe, Liebe, Freundschaft seine Ansprüche so auf Messers Schneide stellen, dass jedes Zuviel oder Zuwenig eine unheilbare Verwundung, eine endgültige Störung des Gleichgewichts bedeutet? Antwort: das darf allenfalls der grosse schöpferische Mensch, dessen capriciöse Neigungen, wie bei schwangern Frauen, weniger ihm selbst als dem werdenden Neuen zur Last fallen — aber nimmermehr der Alltagsmann und die Alltagsfrau der modernen Dichtung, Wesen, deren ganze Productivität in der liebevollen Grosszüchtung ihrer Sonderlings-Unarten und nervösen Monomanien besteht. — An dieser Stelle aber schüttelt die „Seelendichtung“ unsere Vorwürfe hohnlächelnd ab, denn ihr ist der Alltag die Grotte der ungehobenen Schätze; sie hat ihn zu Ehren und Göttlichkeit gebracht. „Wäre ich Plato, Pascal oder Michelangelo, und meine Geliebte schwatzte mir von Ohrringen, — meine Worte wie ihre Worte schwämmen als der gleiche Schaum auf der Oberfläche des Innenmeeres. Mein erhabenster Gedanke wiegt nicht mehr in den Schalen des Lebens und der Liebe als die paar dummen Worte, worin mir ein verliebtes Mädel von seinem Silberring oder Halsband erzählt \Dots{3} Sei noch so gross, weise, redebegabt: die Seele des Bettlers, der im Winkel kauert, beneidet dich nicht, aber deine Seele beneidet vielleicht ihn um sein Schweigen. Der Held bedarf der Anerkennung des Alltagsmenschen, aber dieser verlangt nicht die Anerkennung des Helden; ruhig geht er seinen Gang wie Einer, der seine Schätze sicher geborgen weiss.“ Diese stupenden Sätze Maeterlincks zeigen, dass auch das neue Seelenevangelium ein \Foreign{évangile des humbles} und ein Sclavenaufstand zu werden droht: und die Ersten werden die Letzten sein \Dots{4}
\par

